\section{Problem statement}\label{sec:problem}
Evaluating a synthetic LiDAR dataset reliably requires training a perception model on it and testing on held-out real data. This is expensive, because it must be repeated for every candidate dataset or sensor setting. Cheaper data-level diagnostics such as FRD, BEV MMD and JSD are therefore widely used to judge generated scans~\cite{zyrianovLearningGenerateRealistic2022}. These scores summarise global feature or density distributions. The reviewed evidence suggests that they can disagree with each other~\cite{zyrianovLearningGenerateRealistic2022}, that no single score establishes training utility, and that aggregate scores need not track downstream outcomes~\cite{manivasagamZeroDomainGap2023}. Matching common observations may also hide the rare, small or distant cases that synthesis is meant to add.

It is plausible that such diagnostics catch large failures but miss subtle ones. To our knowledge, this has not been quantified: it is unclear how much a sensor property must degrade before the diagnostics register the change, compared with the point where class- and range-specific real-data performance drops. Without this, practitioners cannot tell when a good diagnostic score is a safe reason to skip real-data evaluation, and when it is not.
% TODO: back the novelty claim with the documented search (Scopus/Scholar) from the review.